\section{Introduction}
Understanding a painting is not a single act of perception. It is a layered process that draws simultaneously on visual attention, affective response, formal knowledge, and historical memory~\cite{BMVC_style_2013,panofsky1955meaning, baxandall1988painting}. A viewer encountering Vermeer's Girl with a Pearl Earring may first register the luminous skin and direct gaze — a narrative observation. A trained eye may then notice the sfumato technique and restricted palette — a formal judgment. The same viewer may feel unease or intimacy — an emotional response. And an art historian will situate the work within the Dutch Golden Age domestic portrait tradition — a contextual interpretation. These four acts of understanding are not sequential stages but concurrent, interacting perspectives on a single visual object.

Current vision-language models do not work this way. Despite rapid advances in multimodal understanding~\cite{NEURIPS2023_llava, chen2023sharegpt4v, qwen2vl, yu2025crispsam2sam2crossmodalinteraction}. , existing systems are predominantly trained on object-centric, single-perspective captions that flatten this interpretive complexity into a single descriptive register. Benchmarks and datasets in computational art analysis — SemArt~\cite{semart_2018}, ExpArt~\cite{ExpArt_acl2024}, ArtEmis~\cite{achlioptas2021artemis}, OmniArt~\cite{2017omniart} — have each captured important dimensions of art understanding, but none provides a unified, large-scale resource that systematically represents all four interpretive dimensions simultaneously for the same artworks. As a result, models trained on these resources learn to describe paintings, but not to interpret them in the full sense that word carries in art history and criticism.

We introduce MMArt, a large-scale multi-perspective dataset of over 74k paintings annotated across four interpretive dimensions: Narrative and Scene Interpretation, Formal Visual Analysis, Emotional Response, and Historical and Contextual Analysis. Each painting in MMArt is paired with four independently generated captions, one per perspective, produced through a hybrid annotation pipeline combining state-of-the-art multimodal models including Qwen3-VL~\cite{bai2025qwen3vltechnicalreport}, GalleryGPT~\cite{MM24GalleryGPT}, ArtRAG~\cite{ArtRAG_MM2025}, and ArtEmis~\cite{achlioptas2021artemis}. The perspectives are generated independently to preserve interpretive diversity and merged into a unified polysemic caption through an explanation harmonization step. MMArt is the first dataset to provide this breadth of interpretive coverage at this scale, with open licensing and reproducible generation protocols.


We demonstrate this through two experimental contributions that extend the dataset into an active benchmark. First, we introduce a perspective sufficiency analysis using painting regeneration: given only the textual perspectives of a painting — without the original image — we test whether generative models can reconstruct the visual content, and we measure which perspective subsets are sufficient to recover which dimensions of the original (style, composition, iconographic content). This provides an empirical grounding for the claim that MMArt's perspectives encode distinct and complementary information rather than paraphrasing each other. Second, we conduct Multi-Perspective Cross-Modal Retrieval. We show that unified multi-perspective captions consistently outperform any single-perspective query on standard retrieval metrics (Recall@1, Recall@5, Recall@10), providing direct evidence that the four perspectives encode complementary information that jointly produces richer visual grounding than any individual interpretive dimension alone.



\noindent Our contributions are:
\begin{itemize}
    \item \textbf{Multi-Perspective Dataset.} We introduce MMArt, the first large-scale multimodal dataset providing four independently generated interpretive perspectives for approximately 70,000 paintings, covering narrative, formal, emotional, and historical dimensions simultaneously for the same artworks. Each painting additionally receives a harmonized unified caption synthesizing all four perspectives.

    \item \textbf{Perspective Complementarity Benchmark.} We operationalize the complementarity of perspectives through an image reconstruction experiment: given text from one or more perspective subsets, a generative model attempts to reconstruct the original painting, and fidelity is measured across style, composition, iconographic content, and emotional tone. This provides a falsifiable empirical test of whether the perspectives encode distinct visual information.
\end{itemize}